\section{Support for Existing Benchmarks}
\label{app:benchmark-suites}

In addition to circuits generated by \sys, the evaluation pipeline can ingest circuits
from established benchmark suites: SupermarQ~\cite{tomesh2022supermarqscalablequantumbenchmark},
MQT Bench~\cite{Quetschlich_2023}, and QASM
Bench~\cite{li2022qasmbenchlowlevelqasmbenchmark}. These suites provide workloads that complement our generated corpus. The compositional circuit simulation workload introduced in Sec.~\ref{sec:benchmark}
remains necessary for controlled simulation scale, feature-space coverage, and reproducible
randomized workload generation, while external suites provide an additional validation
source based on named benchmark circuits.

\begin{table}[t]
  \centering
  \caption{External benchmark-suite circuit entries supported by the \sys ingestion pipeline. Per-suite counts are before cross-suite deduplication.}
  \label{tab:benchmark-suite-ingestion}
  \footnotesize
  \begin{tabular}{@{}lr@{}}
    \toprule
    \textbf{Benchmark suite} & \textbf{\# circuit entries} \\
    \midrule
    SupermarQ & 8 \\
    MQT Bench & 34 \\
    QASM Bench & 67 \\
    \midrule
    Total (after deduplication) & 85 \\
    \bottomrule
  \end{tabular}
\end{table}

All imported circuits are normalized to Qiskit \texttt{QuantumCircuit} objects and
tagged with their source suite and benchmark name. The ingestion pipeline then applies
the same SQL generation, feature extraction, and the rest of the workflow used for \sys-generated
circuits. As a result, external-suite circuits can be queried, filtered, and evaluated
together with generated circuits under the same RDBMS, simulator-selection, and
out-of-core workflows.

The per-benchmark counts in Table~\ref{tab:benchmark-suite-ingestion} are before deduplication. Since
the three benchmark suites contain overlapping algorithm families and problem-size variants, \sys
also records a deduplicated inventory; after deduplication, the external benchmark
inventory covers 85 unique algorithm families. The complete per-suite circuit-family
list and ingestion commands are provided in the \emph{Benchmark Algorithms} table in the online documentation.\footnote{\url{https://github.com/InfiniData-Lab/InferQ/blob/main/scripts/benchmark_suites/README.md}}

 
% \section{Deep Profiling of SQLite and Qiskit Aer}
% \label{sec:deep-profile-sqlite-aer}

% % The aggregate evaluation shows that SQLite can be competitive with Qiskit Aer, especially for peak memory, but the aggregate win rates do not explain where this advantage comes from. We have generated and simulated 106 circuits, and profiled the SQLite and Qiskit Aer\footnote{These 106 circuits are sampled from the \num{7705} circuits used in the Sec. 6.}. The goal is to identify the mechanism behind the crossover: when SQLite is faster, whether this is due to sparse relational representation and low overhead; and when Aer is faster, whether SQLite is dominated by SQL execution-plan cost.

% To understand why SQLite uses less peak memory than
% Qiskit Aer on many circuits, we have generated and simulated 106 circuits, and profiled the SQLite and Qiskit Aer\footnote{These 106 circuits are sampled from the \num{7705} circuits used in the Sec. 6.}.
% % This analysis separates SQLite's memory advantage from possible simulator memory
% % overheads and identifies when SQLite's cost is dominated by SQL execution.

% For SQLite, each circuit is compiled into a tensor-contraction program represented as a chain of Common Table Expressions (CTEs). Each contraction step is implemented as a join followed by an aggregation. We record the number of total and contraction CTEs, estimated intermediate-relation sizes, result cardinalities, SQLite execution-plan operators, temporary B-tree usage, automatic-index usage, page-cache overflow counters, peak resident set size (RSS), and wall-clock runtime. For Qiskit Aer, we compare against the fastest successful Aer method for the same circuit, rather than fixing a single simulator representation across all circuits.

% The profiling results in Table~\ref{tab:sqlite-aer-profile} help us understand \emph{when} each engine wins. Across 106 circuits, SQLite is faster in 19 cases, while Qiskit Aer is faster in 87 cases. SQLite wins when the relational workload remains small: the median SQLite-winning circuit has 37 total CTEs, 32 contraction CTEs, a largest intermediate relation of 256\,B, and only 16 final result rows. In these cases, SQLite runs in a median of 0.85\,ms, compared with 2.01\,ms for the fastest successful Aer method. By contrast, when Aer wins, the median workload is larger: 55 total CTEs, 49 contraction CTEs, a largest intermediate relation of 1.0\,KB, and 128 final result rows. In these cases, SQLite takes a median of 1.67\,ms, while Aer takes 0.92\,ms.

% \begin{table}[t]
% \centering
% \small
% \caption{Median properties of completed SQLite--Aer profiling runs, grouped by the faster runtime backend. CTE sizes are estimated from intermediate-relation cardinalities and row widths.}
% \label{tab:sqlite-aer-profile}
% \begin{tabular}{lrr}
% \toprule
%  & SQLite faster & Aer faster \\
% \midrule
% Number of cases & 19 & 87 \\
% Total CTEs & 37 & 55 \\
% Contraction CTEs & 32 & 49 \\
% Total estimated CTE size & 3.0\,KB & 8.6\,KB \\
% Largest intermediate CTE & 256\,B & 1.0\,KB \\
% Final result rows & 16 & 128 \\
% SQLite runtime & 0.85\,ms & 1.67\,ms \\
% Fastest Aer runtime & 2.01\,ms & 0.92\,ms \\
% \bottomrule
% \end{tabular}
% \end{table}

% This result shows that SQLite's advantage is not determined by qubit count alone. Some SQLite-winning circuits have more qubits than Aer-winning circuits, but their relational intermediates remain small. The decisive factor is the size and structure of the induced SQL workload. When the tensor-contraction sequence preserves sparsity, SQLite only materializes a small number of non-zero amplitudes, and the join-and-aggregate pipeline can finish before Aer method setup and simulation overhead dominate. For such circuits, the sparse SQL representation is a good representation of the quantum state in the circuits.

% The same representation becomes unfavorable once the contraction pipeline grows. Longer CTE chains introduce repeated joins, aggregations, and materialization points. SQLite query plans repeatedly use temporary B-trees for grouping and automatic indexes for joins. These are standard relational execution mechanisms, but in a long contraction sequence their cumulative cost dominates the arithmetic. Thus, when Aer wins, it is not because SQLite fails to represent the computation; rather, the relational execution machinery becomes the bottleneck.

% % The timeout cases follow the same pattern. They occur in the high-cost SQL regime, with larger join-and-aggregation structure: the median timeout case has 15 qubits, 99 \texttt{GROUP BY} clauses, and 128 join/equality predicates. These cases indicate that the hard boundary for SQLite is not simply state-vector dimension, but the interaction between output density, intermediate cardinality, CTE-chain length, and the physical operators selected by SQLite. Dense or highly connected circuits can rapidly turn the SQL representation into a long sequence of expensive relational operators.

% The memory measurement results are different from the runtime results. Among 106 circuits, SQLite uses less peak RSS than the best Aer memory configuration in 105 out of 106 cases. Median peak RSS is 167\,MB for SQLite and 171\,MB for Aer. This difference is small in absolute terms for these circuits, but it is consistent: SQLite benefits from storing sparse relational intermediates rather than allocating the simulator representation chosen by Aer. At the same time, this memory advantage does not automatically translate into runtime advantage, because SQLite still pays the cost of repeated SQL planning and execution operators.
% %
% The SQLite profiling results are similar. The page-cache overflow counter is zero in all completed runs, and we do not observe an external-sort. Therefore, the slowdown is not explained by SQLite spilling large page-cache contents to disk. Instead, it might be because the repeated use of temporary B-trees and automatic indexes inside the query plan. 
% % In this setting, temporary B-trees should be interpreted primarily as relational execution structures used for grouping and ordering, not as evidence of out-of-core execution. 
% The dominant cost is the repeated construction and consumption of these structures across the CTE chain.

% Different from RDBMSs, Aer's performance depends strongly on which simulator method is valid and efficient for the circuit. In the profiling runs, the fastest successful Aer method varies across circuits: \texttt{unitary} is fastest in 47 cases, \texttt{automatic} in 40 cases, \texttt{extended\_stabilizer} in 12 cases, and \texttt{statevector} in 7 cases. This confirms that a single fixed Aer method is not a stable baseline for heterogeneous circuits. Some circuits admit efficient specialized simulation, while others require a more general representation. For peak memory, the best Aer configuration is almost always \texttt{statevector}, which is the lowest-RSS Aer method in 105 completed cases.

% Overall, the profiling study explains the SQLite--Aer crossover. SQLite wins when three conditions hold simultaneously: the circuit state remains sparse, the intermediate relations stay small, and the CTE chain is short enough that SQL execution overhead is negligible. Aer wins once the contraction sequence becomes long enough for repeated joins, aggregations, temporary B-trees, and automatic indexes to dominate. The comparison is therefore governed by both data representation and execution strategy: SQLite provides a compact sparse representation for some circuits, while Aer benefits from specialized simulator methods that avoid the overhead of executing a long relational contraction pipeline.

\section{Profiling of SQLite and Qiskit Aer}
\label{sec:deep-profile-sqlite-aer}

To understand why SQLite uses less peak memory than Qiskit Aer on many circuits,
we profile both systems on 162 generated circuits.\footnote{We have sampled these 162 circuits
  from the large \num{7705} circuits used in Section~6. This profiling comparison is separate from the all-backend tuned
winner count reported in Figure~7 and Appendix~G.}

For SQLite, each circuit is compiled into a tensor-contraction program represented
as a chain of Common Table Expressions (CTEs). Each contraction step is implemented
as a join followed by an aggregation. We record the number of total and contraction
CTEs, estimated intermediate-relation sizes, result cardinalities, SQLite
execution-plan operators, temporary B-tree usage, automatic-index usage,
page-cache overflow counters, peak resident set size (RSS), and wall-clock
runtime. For Qiskit Aer, we compare against the fastest   Aer method for
runtime, rather than fixing a single simulator representation across all circuits.
For memory, we compare against the lowest-RSS   Aer configuration for
the same circuit.

The profiling results\footnote{\url{https://github.com/InfiniData-Lab/InferQ/blob/main/analysis/memory_profiling/sqlite_vs_qiskit_summary.csv}} in Table~\ref{tab:sqlite-aer-profile} characterize when
each system is faster. Across 162 circuits, SQLite is faster in 46 cases, while
Qiskit Aer is faster in 116 cases. SQLite is faster when the relational workload
remains small: the median SQLite-winning circuit has 26.5 total CTEs, 22
contraction CTEs, the largest intermediate relation of 256\,B, and 16 final result
rows. In these cases, SQLite runs in a median of 2.03\,ms, compared with
2.82\,ms for the fastest   Aer method. In contrast, when Aer is faster,
the median workload is larger: 55 total CTEs, 49 contraction CTEs, a largest
intermediate relation of 1.0\,KB, and 128 final result rows. In these cases,
SQLite takes a median of 6.63\,ms, while Aer takes 3.09\,ms.

\begin{table}[t]
\centering
\small
\caption{Median properties of completed SQLite--Aer profiling runs, grouped by
the backend with lower runtime. CTE sizes are estimated from intermediate-relation
cardinalities and row widths.}
\label{tab:sqlite-aer-profile}
\begin{tabular}{lrr}
\toprule
 & SQLite faster & Aer faster \\
\midrule
Number of cases & 46 & 116 \\
Total CTEs & 26.5 & 55 \\
Contraction CTEs & 22 & 49 \\
Total estimated CTE size & 2.9\,KB & 10.5\,KB \\
Largest intermediate CTE & 256\,B & 1.0\,KB \\
Final result rows & 16 & 128 \\
SQLite runtime & 2.03\,ms & 6.63\,ms \\
Fastest Aer runtime & 2.82\,ms & 3.09\,ms \\
\bottomrule
\end{tabular}
\end{table}

These results show that SQLite's advantage is not determined by qubit count
alone. Some SQLite-winning circuits have more qubits than Aer-winning circuits,
but their relational intermediates remain small. The decisive factor is the size
and structure of the induced SQL workload. When the tensor-contraction sequence
preserves sparsity, SQLite materializes only a small number of nonzero amplitudes,
and the join-and-aggregate pipeline can complete before Aer's method setup and
simulation overheads dominate. For such circuits, the sparse SQL representation
provides a compact representation of the circuit state.

The same representation becomes less favorable once the contraction pipeline
grows. Longer CTE chains introduce repeated joins, aggregations, and intermediate
execution structures. SQLite query plans use temporary B-trees for grouping in
all completed runs and automatic indexes for joins in 154 out of 162 runs. These
are standard relational execution mechanisms, but in a long contraction sequence
their cumulative cost can dominate the arithmetic. Thus, when Aer is faster, it
is not because SQLite fails to represent the computation; rather, the relational
execution machinery becomes the bottleneck.

The memory measurements follow a different pattern from the runtime measurements.
Among the 162 circuits, SQLite uses less peak RSS than the lowest-RSS  
Aer method in 139 cases. Median peak RSS is 203\,MB for SQLite and 206\,MB for
Aer. This difference is small in absolute terms for these circuits, but it is
consistent with SQLite benefiting from storing sparse relational intermediates
rather than allocating the simulator representation used by Aer. At the same
time, this memory advantage does not automatically translate into a runtime
advantage, because SQLite still pays the cost of repeated SQL execution.

The SQLite-specific profiling counters support this interpretation. The
page-cache overflow counter is zero in all completed runs, and we do not observe
external-sort behavior. Thus, the slowdown is not explained by observed
out-of-core execution in these runs. Instead, the dominant cost is the repeated
construction and consumption of temporary B-trees and automatic indexes across
the CTE chain.

Unlike the RDBMS backends, Aer's performance depends strongly on which simulator
method is applicable and efficient for each circuit. In the profiling runs, the
fastest  Aer method varies across circuits: \texttt{unitary} is fastest
in 46 cases, \texttt{automatic} in 42 cases, \feat{statevector} in 37 cases,
\feat{extended\_stabilizer} in 26 cases, \feat{matrix\_product\_state} in
5 cases, \feat{density\_matrix} in 4 cases, and \feat{stabilizer} in 2
cases. This indicates that a single fixed Aer method is not a stable baseline for
heterogeneous circuits. Some circuits admit efficient specialized simulation,
while others require a more general representation. For peak memory, the
lowest-RSS Aer configuration is usually \feat{statevector} or
\feat{density\_matrix}:  \feat{statevector} has the lowest RSS in 88 cases,
\feat{density\_matrix} in 49 cases, \feat{automatic} in 14 cases,
\feat{matrix\_product\_state} in 9 cases, and \feat{unitary} in 2 cases.

We do not enable Aer cache blocking in this comparison. Our Aer baseline uses the
standard CPU simulator methods, whereas Aer's cache-blocking options (e.g.,
\texttt{blocking\_enable} and \texttt{blocking\_qubits}) are documented for
distributed GPU and/or multi-node simulation~\cite{QiskitAerMultiGPU}. They
partition the state into chunks to reduce data exchange across distributed memory
spaces, rather than providing a disk-backed CPU cache comparable to an RDBMS
buffer manager. Enabling these options would therefore change the simulator
configuration instead of providing a fair comparison to SQLite's
relational execution.  

\medskip
  \vspace{0.2cm}
\noindent
\setlength{\fboxsep}{6pt}
\fcolorbox{black!15}{black!4}{%
  \parbox{\dimexpr\linewidth-2\fboxsep-2\fboxrule\relax}{
\noindent\textbf{Takeaway.}
The relative performance of SQLite and Qiskit Aer is governed by both
representation and execution strategy. SQLite is memory-efficient when sparsity
keeps intermediate relations small, but Aer is faster once repeated joins,
aggregations, temporary B-trees, and automatic indexes make SQL execution
overhead dominate the contraction pipeline.
}} 

% \section{Additional benchmark-suite circuits}
% \label{app:benchmark-suites}

% In addition to the randomly generated circuits used by \sys, the evaluation
% pipeline can ingest circuits from established quantum benchmark suites:
% SupermarQ~\cite{tomesh2022supermarqscalablequantumbenchmark}, MQT
% Bench~\cite{Quetschlich_2023}, and QASM
% Bench~\cite{li2022qasmbenchlowlevelqasmbenchmark}. These suites provide
% externally curated circuits that complement the synthetic generator. The random
% generator is useful for broad, controlled coverage of the feature space, while
% benchmark-suite circuits provide recognizable algorithmic workloads for
% cross-checking that trends observed on generated circuits also appear on
% community benchmarks. After deduplicating equivalent algorithm families across
% suites and problem sizes, the benchmark-suite inventory covers 85 unique
% quantum algorithm families.

% InferQ includes standalone loaders under
% \path{scripts/benchmark_suites}. Each loader produces Qiskit
% \texttt{QuantumCircuit} objects and tags each circuit with a source-suite
% identifier and a benchmark name. The ingestion script
% \path{scripts/ingest_benchmarks.py} then passes these circuits through the
% same feature extraction, simulation, duplicate detection, hashing, and local
% storage path used for generated InferQ circuits. As a result, benchmark-suite
% circuits can be merged into the same metadata tables and evaluated using the
% same RDBMS, simulator-selection, and out-of-core workflows.

% \begin{table}[t]
%   \centering
%   \caption{External benchmark-suite support in the InferQ ingestion pipeline.}
%   \label{tab:benchmark-suite-ingestion}
%   \footnotesize
%   \begin{tabularx}{\columnwidth}{@{}lXX@{}}
%     \toprule
%     \textbf{Suite} & \textbf{Quantum Algorithms} & \textbf{Ingestion notes} \\
%     \midrule
%     SupermarQ & 8 benchmark entries & Converted from Cirq to Qiskit \\
%     MQT Bench & 34 benchmark entries & Loaded through the MQT Bench API \\
%     QASM Bench & 67 benchmark entries & Parsed from a pinned QASM snapshot \\
%     \bottomrule
%   \end{tabularx}
% \end{table}

% For SupermarQ, the loader covers the application benchmarks exposed by the
% suite, including GHZ, Hamiltonian simulation, Mermin-Bell, error-correction,
% QAOA, and VQE-style circuits. For MQT Bench, the loader uses a curated
% algorithm-level subset spanning state preparation, Fourier and phase-estimation
% circuits, oracle algorithms, variational circuits, quantum walks, random
% circuits, and representative arithmetic/error-correction circuits. For QASM
% Bench, InferQ vendors the small and medium tiers as OpenQASM files, records the
% upstream commit in a manifest, and excludes the large tier to keep the
% evaluation within the same qubit and runtime budgets as the rest of the paper.

% This design makes the external suites an additive evaluation source rather than
% a replacement for the generated corpus. Generated circuits remain necessary for
% systematic coverage, controllable scale, and reproducibility under fixed random
% seeds. The external suites provide an additional validity check by exercising
% InferQ on named circuits that are already used by the quantum-computing
% community.



% \begin{table*}[t]
%   \centering
%   \caption{Two-profile database-specific tuning settings used for the 7{,}705-circuit evaluation.}
%   \label{tab:finetuned-rdbms-profiles}
%   \footnotesize
%   \setlength{\tabcolsep}{4pt}
%   \begin{tabularx}{\textwidth}{@{}llXX@{}}
%     \toprule
%     \textbf{Engine} & \textbf{Profile} & \textbf{Parameters} & \textbf{Rationale} \\
%     \midrule
%     DuckDB
%       & \texttt{small\_size}
%       & \texttt{memory\_limit=10780MB}; \texttt{threads=2};
%         \texttt{preserve\_insertion\_order=false};
%         \texttt{max\_temp\_directory\_size=256GB}
%       & Gives DuckDB enough buffer-pool memory and limited parallelism for
%         short queries, while disabling insertion-order preservation to allow
%         spill-capable operators. \\
%     DuckDB
%       & \texttt{large\_size}
%       & \texttt{memory\_limit=7409MB}; \texttt{threads=1};
%         \texttt{preserve\_insertion\_order=true};
%         \texttt{max\_temp\_directory\_size=128GB}
%       & Reduces concurrent memory pressure and temp-space exposure for larger
%         query plans, where planner/runtime overhead can dominate. \\
%     SQLite
%       & \texttt{small\_size}
%       & \texttt{cache\_mb=64}; \texttt{db\_path=:memory:};
%         \texttt{mmap\_mb=0}; \texttt{temp\_store=MEMORY};
%         \texttt{threads=1}
%       & Keeps small workloads fully in memory and avoids filesystem overhead. \\
%     SQLite
%       & \texttt{large\_size}
%       & \texttt{cache\_mb=128}; file-backed database;
%         \texttt{mmap\_mb=128}; \texttt{temp\_store=MEMORY};
%         \texttt{threads=1}
%       & Uses a persistent database file and modest mmap window to stabilize
%         larger intermediate results while retaining memory-backed temporaries. \\
%     PostgreSQL
%       & \texttt{small\_size}
%       & \texttt{work\_mem=256MB}; \texttt{temp\_buffers=128MB};
%         \texttt{effective\_cache\_size=8GB};
%         \texttt{hash\_mem\_multiplier=2.0};
%         \texttt{max\_parallel\_workers\_per\_gather=4};
%         \texttt{join\_collapse\_limit=1}; \texttt{from\_collapse\_limit=1};
%         \texttt{jit=off}
%       & Allocates more per-query memory and parallelism to accelerate small
%         contractions, while preserving the emitted join order and avoiding JIT
%         overhead. \\
%     PostgreSQL
%       & \texttt{large\_size}
%       & \texttt{work\_mem=55MB}; \texttt{temp\_buffers=65MB};
%         \texttt{effective\_cache\_size=2GB};
%         \texttt{hash\_mem\_multiplier=1.94};
%         \texttt{max\_parallel\_workers\_per\_gather=0};
%         \texttt{join\_collapse\_limit=4}; \texttt{from\_collapse\_limit=1};
%         \texttt{jit=off}
%       & Caps per-node memory and disables parallel workers for large plans,
%         reducing worker amplification and making externalization behavior more
%         predictable. \\
%     Umbra
%       & comparison engine
%       & PostgreSQL-compatible server interface; reported version:
%         \texttt{PostgreSQL 18.1 compatible umbra v0.2-1665-gfeeb4bb5d}
%       & Broadens the engine comparison with a modern compiled analytical DBMS
%         that accepts the same SQL workload through a PostgreSQL-compatible
%         client path. \\
%     \bottomrule
%   \end{tabularx}
% \end{table*}
\begin{table*}[t]
  \centering
  \caption{Database-specific tuning settings used for the 7{,}705-circuit evaluation. 
  % The rationale for each profile is discussed in Section~\ref{app:mlos-runner}.
  }
  \label{tab:finetuned-rdbms-profiles}
  \footnotesize
  \setlength{\tabcolsep}{4pt}
  \begin{tabularx}{\textwidth}{@{}llX@{}}
    \toprule
    \textbf{Engine} & \textbf{Profile} & \textbf{Parameters} \\
    \midrule
    DuckDB
      & \texttt{small\_size}
      & \texttt{memory\_limit=10780MB}; \texttt{threads=2};
        \texttt{preserve\_insertion\_order=false};
        \texttt{max\_temp\_directory\_size=256GB} \\
    DuckDB
      & \texttt{large\_size}
      & \texttt{memory\_limit=7409MB}; \texttt{threads=1};
        \texttt{preserve\_insertion\_order=true};
        \texttt{max\_temp\_directory\_size=128GB} \\
    SQLite
      & \texttt{small\_size}
      & \texttt{cache\_mb=64}; \texttt{db\_path=:memory:};
        \texttt{mmap\_mb=0}; \texttt{temp\_store=MEMORY};
        \texttt{threads=1} \\
    SQLite
      & \texttt{large\_size}
      & \texttt{cache\_mb=128}; file-backed database;
        \texttt{mmap\_mb=128}; \texttt{temp\_store=MEMORY};
        \texttt{threads=1} \\
    PostgreSQL
      & \texttt{small\_size}
      & \texttt{work\_mem=256MB}; \texttt{temp\_buffers=128MB};
        \texttt{effective\_cache\_size=8GB};
        \texttt{hash\_mem\_multiplier=2.0};
        \texttt{max\_parallel\_workers\_per\_gather=4};
        \texttt{join\_collapse\_limit=1}; \texttt{from\_collapse\_limit=1};
        \texttt{jit=off} \\
    PostgreSQL
      & \texttt{large\_size}
      & \texttt{work\_mem=55MB}; \texttt{temp\_buffers=65MB};
        \texttt{effective\_cache\_size=2GB};
        \texttt{hash\_mem\_multiplier=1.94};
        \texttt{max\_parallel\_workers\_per\_gather=0};
        \texttt{join\_collapse\_limit=4}; \texttt{from\_collapse\_limit=1};
        \texttt{jit=off} \\
    % Umbra
    %   & comparison engine
    %   & PostgreSQL-compatible server interface; reported version:
    %     \texttt{PostgreSQL 18.1 compatible umbra v0.2-1665-gfeeb4bb5d} \\
    \bottomrule
  \end{tabularx}
\end{table*}


% \section{DBMS-Specific Tuning }
% \label{app:finetuned-rdbms}

% To understand how RDBMS configuration affects simulation performance, we have added  
% DBMS-specific tuning for PostgreSQL, DuckDB, and SQLite. 
% % We keep the input circuit set and \sys’s SQL generation
% % pipeline unchanged, and introduce a new runner that varies (i) \emph{DBMS configuration
% % parameters} (memory, parallelism, temporary I/O, and optimizer/planner settings) and
% % (ii) \emph{execution-structure choices} exposed by the SQL formulation (CTE
% % materialization and split execution). 
% The code and documentation are updated in the
% anonymous repository.\footnote{\url{}}.


% \medskip
% \para{Automated tuning via configuration search}
% We initially explored manual tuning to understand the impact of major parameters, but the
% configuration space is large, engine-specific, and not scalable across thousands of
% circuits. We therefore adopt an automated configuration-search workflow based on MLOS
% (Machine Learning for Operating Systems)\footnote{\url{https://github.com/microsoft/MLOS}}:
% the workload is fixed, and this tuner searches over engine parameters to minimize runtime.
% % The all-engine tuning experiment uses a 162-circuit validation subset, while the larger
% % SQLite-vs.-Qiskit comparison is run on the 7{,}705-circuit dataset used in
% % Figure~\ref{fig:aer-vs-sqlite-wins}. We keep these two roles separate: the 162-circuit
% % experiment selects and validates DBMS configurations, whereas the 7{,}705-circuit SQLite
% % run evaluates the main memory/runtime comparison at scale. For reproducibility, we
% % document the parameter spaces, seed profiles, and runner details in
% % Section~\ref{app:tuning-repro}.

% \medskip
% \para{Our workload-specific tuning}
% In preliminary experiments, we observed that a \emph{single} global configuration is suboptimal
% because circuits induce heterogeneous query shapes (e.g., very short vs.\ long CTE chains)
% and different bottlenecks (overhead-dominated vs.\ memory/IO-dominated). We therefore
% introduce a workload-specific separation. 
% for example 
% for the 162-circuit all-engine MLOS validation we divide circuits by median
% \texttt{circuit\_size} (33 gates): circuits with \texttt{circuit\_size} $\le 33$ use a
% \texttt{small\_size} profile (82 circuits), and the remaining circuits use a
% \texttt{large\_size} profile (80 circuits).
% This
% allows us to tune configurations to structurally different workloads: 
%  small circuits benefit from
%  settings like larger caches and limited parallelism, while larger
% circuits often require certain settings to reduce planner overhead and avoid
% unpredictable spilling.
% % For example, for the 7{,}705-circuit dataset\footnote{\url{https://anonymous.4open.science/r/InferQ-ADEF/analysis/training_data/rdbms_training_data.parquet}}, we use a separate
% % static-contraction SQLite profile: both size classes keep the database in memory
% % (\texttt{db\_path=:memory:}), use a 64~MB SQLite cache, disable cache spilling, and place
% % temporary structures in file-backed temp storage (\texttt{temp\_store=FILE}). 
% This run is
% not an additional MLOS search; it is the scaled SQLite validation used to compare against
% the best Qiskit Aer method. Table~\ref{tab:finetuned-rdbms-profiles} shows the dominating parameters found by MLOS after we set such a work-specific tuning. 
% % small circuits can benefit from more aggressive settings (e.g., caching), whereas
% % larger circuits often require choices to avoid planner overhead,
% % unfavorable materialization, and unstable spilling behavior.


% % In our SQL-based simulator, the tensor-contraction sequence is written as a chain of Common Table Expressions (CTEs), each encoding one contraction step (join + aggregate). Our preliminary finding is that data representation, SQL execution-plan overhead, and Aer method choice drive the observed differences.
% % We report CTE materialization because it changes how the optimizer/executor
% % treats the SQL query for each circuit (\texttt{WITH ... SELECT}) without converting the
% % workload into a sequence of explicit intermediate relations. In the out-of-core QFT
% % (3--26 qubits) experiment at a 4~GB cap, \texttt{monolithic\_materialized} preserves the
% % same completion range as \texttt{monolithic} for DuckDB (3--24 qubits) and for
% % SQLite/PostgreSQL (3--26 qubits), while reducing matched median runtime relative to
% % \texttt{monolithic} by roughly $2.1\times$ for DuckDB and $1.8\times$ for PostgreSQL;
% % SQLite is essentially unchanged. \texttt{split} remains the most robust out-of-core
% % structure, completing all PostgreSQL and SQLite configurations and all DuckDB
% % configurations except the 26-qubit/4~GB case.

% \medskip
 





% %https://github.com/microsoft/MLOS%

% \begin{table*}[t]
%   \centering
%   \caption{Fine-tuned all-engine validation on the 162-circuit MLOS subset. Each engine is run twice per circuit. Ratios compare the tuned engine to the best successful Qiskit Aer method; values below $1.0\times$ favor the RDBMS engine.}
%   \label{tab:finetuned-rdbms-results}
%   \small
%   \setlength{\tabcolsep}{4pt}
%   \begin{tabular}{lrrrrrrr}
%     \toprule
%     \textbf{Engine} & \textbf{Succ. circuits} & \textbf{Succ. runs} & \textbf{Timeout/error runs} & \textbf{Time wins} & \textbf{Median time} & \textbf{Memory wins} & \textbf{Median memory} \\
%     \midrule
%     SQLite & 136/162 & 270/324 & 54/324 & 7/136 & $4.99\times$ & 95/136 & $0.51\times$ \\
%     PostgreSQL & 134/162 & 267/324 & 57/324 & 0/134 & $12.00\times$ & 9/134 & $2.23\times$ \\
%     DuckDB & 71/162 & 139/324 & 185/324 & 0/71 & $33.68\times$ & 11/71 & $1.35\times$ \\
%     \bottomrule
%   \end{tabular}
% \end{table*}

% \begin{table}[t]
%   \centering
%   \caption{SQLite static-contraction validation on the 7{,}705-circuit dataset. Winner counts compare tuned SQLite, untuned baseline SQLite, and the best successful Qiskit Aer method.}
%   \label{tab:sqlite-7705-static-results}
%   \footnotesize
%   \setlength{\tabcolsep}{3pt}
%   \begin{tabular}{lrrr}
%     \toprule
%     \textbf{Objective} & \textbf{Tuned SQLite} & \textbf{Baseline SQLite} & \textbf{Qiskit Aer} \\
%     \midrule
%     Runtime & 0 (0.0\%) & 1{,}067 (13.8\%) & 6{,}638 (86.2\%) \\
%     Peak memory & 2{,}085 (27.1\%) & 1{,}818 (23.6\%) & 3{,}802 (49.3\%) \\
%     \bottomrule
%   \end{tabular}
% \end{table}

% \para{Results}
% Table~\ref{tab:finetuned-rdbms-results} shows that simulation results after tuning over the 162 circuit datasets\footnote{\url{https://anonymous.4open.science/r/InferQ-ADEF/analysis/training_data/rdbms_all_methods_training_data.parquet}}. We see that PostgreSQL and DuckDB's performance are significatnly improved by automatic tuning, esepcially for peak memeory, athough   Qiskit Aer remains the fastest backend for
% most general circuits.

% We did not observe SQLite's performance did not have an performance gain from autmoatic tuning, so we have conducted manual tuning. Even so, the runtime imporvement is mild 

% The results on 7{,}705-circuit results for SQLite similar. 
% Tuned SQLite improves over the untuned SQLite baseline in 4{,}946/7{,}404
% matched memory comparisons (66.8\%) and uses less peak memory than the best Qiskit Aer
% method in 3{,}891/7{,}404 matched comparisons (52.6\%). Considering the overall winner
% among tuned SQLite, baseline SQLite, and Qiskit Aer, SQLite is the peak-memory winner on
% 3{,}903/7{,}705 circuits (50.7\%): 2{,}085 wins come from tuned SQLite and 1{,}818 from
% the baseline SQLite configuration. Runtime shows the opposite pattern: Qiskit Aer remains
% the winner on 6{,}638/7{,}705 circuits (86.2\%), while all SQLite runtime wins come from
% the untuned baseline configuration (1{,}067/7{,}705 circuits). Thus, SQL-based simulation
% is not universally faster, but its sparse relational representation can be substantially
% more memory-efficient, making learned routing and DBMS-level physical optimization
% valuable.

% % The profile-level breakdown confirms that larger query plans are the harder regime. In
% % the \texttt{small\_size} profile (82 circuits; 164 runs per engine), SQLite completes
% % 153 runs and times out on 11, PostgreSQL completes 152 and times out on 12, and DuckDB
% % completes 97 and times out on 67. In the \texttt{large\_size} profile (80 circuits; 160
% % runs per engine), SQLite completes 117 runs and times out on 43, PostgreSQL completes
% % 115 runs with 44 timeouts and one error, and DuckDB completes 42 runs and times out on
% % 118. No configurations are skipped in this validation run; each engine is attempted for
% % each circuit under its selected profile. By objective, tuned RDBMSs are most helpful for
% % memory on both profiles: they beat Qiskit Aer on 61/78 completed \texttt{small\_size}
% % comparisons and 36/60 completed \texttt{large\_size} comparisons, but runtime wins occur
% % only in the smaller profile (7/78).

% % Table~\ref{tab:sqlite-7705-static-results} reports the actual SQLite static-contraction
% % results on the 7{,}705-circuit dataset. Tuned SQLite completes 7{,}404/7{,}705 circuits
% % (96.1\%) and times out on the remaining 301 circuits; at the run level, this corresponds
% % to 14{,}800 successful runs and 610 timeouts across two runs per circuit. The completion
% % rate is higher for the \texttt{small\_size} profile (3{,}829/3{,}889 circuits) than for the
% % \texttt{large\_size} profile (3{,}575/3{,}816 circuits), confirming that larger query plans
% % remain the harder regime.
\section{DBMS-Specific Tuning}
\label{app:finetuned-rdbms}

To understand how DBMS configuration affects simulation performance, we have added
DBMS-specific tuning for PostgreSQL, DuckDB, and SQLite.\footnote{We did not include Umbra in the DBMS-specific tuning study because its publicly available documentation for the released Docker image only describes basic access through the PostgreSQL-compatible server and command-line interface, and we could not identify a stable, documented set of engine-level tuning parameters needed to define a reproducible configuration-search space: \url{https://hub.docker.com/r/umbradb/umbra }} We keep the circuit set and
\sys’s SQL generation pipeline unchanged and only vary DBMS-level configuration.

\medskip
\para{Automated tuning via configuration search}
We initially explored manual tuning to understand the impact of major parameters, but the
configuration space is large, engine-specific, and not scalable across thousands of
circuits. We therefore adopt an automated configuration-search workflow based on MLOS
(Machine Learning for Operating Systems)\footnote{\url{https://github.com/microsoft/MLOS}}:
the workload is fixed, and the tuner searches over DBMS parameters to minimize runtime.
For reproducibility, we document the parameter spaces, seed profiles, and runner details in
Section~\ref{app:tuning-repro}.

\medskip
\para{Workload-specific tuning}
In preliminary experiments, we observed that a single global configuration is suboptimal:
circuits induce heterogeneous query shapes (e.g., short vs.\ long CTE chains) and
different bottlenecks (overhead-dominated vs.\ memory- and I/O-dominated). We therefore
introduce workload stratification by circuit size. For the 162-circuit all-engine MLOS
validation set,\footnote{\url{https://github.com/InfiniData-Lab/InferQ/blob/main/analysis/training_data/rdbms_all_methods_training_data.parquet}}
we split circuits by the median \texttt{circuit\_size} (33 gates): circuits with
\texttt{circuit\_size} $\le 33$ use a \texttt{small\_size} profile (82 circuits) and the
remaining circuits use a \texttt{large\_size} profile (80 circuits). This allows the tuner
to select configurations appropriate for structurally different workloads of small and large circuits.
% benefit from settings such as larger caches and limited parallelism, whereas larger
% circuits often require more conservative settings to reduce planner overhead and avoid
% unpredictable spilling. 
Table~\ref{tab:finetuned-rdbms-profiles} summarizes the final
profiles. We observed that SQLite did not obtain much performance gain from automatic tuning, so we have tuned it manually\footnote{Tuning results for the 7705 dataset can be found here: \url{https://github.com/InfiniData-Lab/InferQ/blob/main/analysis/finetuned/rdbms7705_sqlite_static_contr/rdbms7705_sqlite_static_contr_results.csv}}.

\begin{table}[t]
  \centering
  \caption{Tuned DB performance on the 162 circuits. Time and peak-memory wins are counted against the best successful Qiskit Aer method per circuit.}
  \label{tab:finetuned-rdbms-results}
  \small
  \setlength{\tabcolsep}{6pt}
  \begin{tabular}{lrr}
    \toprule
    \textbf{Backend} & \textbf{Time wins} & \textbf{Memory wins} \\
    \midrule
    SQLite      & 34/162  & 108/162 \\
    DuckDB      & 0/162   & 0/162 \\
    PostgreSQL  & 0/162   & 0/162 \\
    Umbra       & 0/162   & 0/162 \\
    \midrule
    RDBMS total     & 34/162  & 108/162 \\
    Qiskit Aer total & 128/162 & 54/162 \\
    \bottomrule
  \end{tabular}
\end{table}

\para{Results}
Table~\ref{tab:finetuned-rdbms-results} reports results after tuning 
% \footnote{Tuni \url{https://anonymous.4open.science/r/InferQ-ADEF/analysis/finetuned/rdbms162_all_engines_120s/rdbms162_all_engines_120s_results.csv}}
on the 162-circuit dataset, which are also the results for Figure~\ref{fig:rdbmsqiskitcomparison}. Overall, Qiskit Aer
remains the runtime winner on most circuits (128/162). Among the evaluated DBMS backends,
only SQLite achieves runtime wins (34/162); DuckDB, PostgreSQL, and Umbra do not win on
runtime on this subset. In contrast, SQLite is the peak-memory winner on 108/162 circuits ($66.7\%$),
while Qiskit Aer wins peak memory on 54/162 circuits.
We have also evaluated the tuned configurations on the full \num{7705}-circuit dataset
(Figure~\ref{fig:aer-vs-sqlite-wins}). The result is similar: RDBMS backends achieve the
lowest runtime on \num{1101} circuits (14.3\%), and all of these runtime wins come from the SQLite, for which tuning did not have much improvement. For peak memory, RDBMS backends win on \num{3902}
circuits (50.6\%), dominated by tuned SQLite (\num{3901} wins) with one additional win
by Umbra. Relative to the untuned setting, tuning increases SQLite's peak-memory win count
from \num{3196} to \num{3901} circuits.

\medskip
  \vspace{0.2cm}
\noindent
\setlength{\fboxsep}{6pt}
\fcolorbox{black!15}{black!4}{%
  \parbox{\dimexpr\linewidth-2\fboxsep-2\fboxrule\relax}{
\para{Take-away}
Taken together with Appendix~A (out-of-core) and Appendix~F (memory profiling), these
results suggest two concrete future optimization opportunities for DBMS-backed simulation. 
\\i) For
\emph{large} circuits where memory is the bottleneck, DBMS execution is uniquely useful
because it can externalize operator workspaces and intermediate relations. For such simulation,
end-to-end performance is largely driven by spill and intermediate growth, motivating
spill-aware planning (join ordering, operator selection, and materialization control) and  physical design studies over
intermediate tables (indexes, partitioning and layout, and caching reused intermediates) as discussed in Appendix~H.
\\ii) For \emph{small} circuits where Aer finishes in milliseconds, DBMS runtime is often
dominated by fixed overheads (SQL compilation, planning, and repeated query execution),
motivating workload-level optimizations such as batching,  reuse, and
handling repeated subcircuit computations across many runs. 
\\
\sys is a convenient tool
for future database research in these directions, because it generates controlled workloads
(small or large, sparse or dense) and emits database-native SQL along with query-shape and
circuit metadata needed for systematic evaluation. 
}}
%   \begin{table}[t]
%     \centering
%     \caption{Tuned DB performance on the 162 circuits. Time and
%   memory wins are compared to the best successful Qiskit Aer method.}
%     \label{tab:finetuned-rdbms-results}
%     \small
%     \setlength{\tabcolsep}{6pt}
%     \begin{tabular}{lrr}
%       \toprule
%       \textbf{Engine} & \textbf{Time wins} & \textbf{Memory wins} \\
%       \midrule
%       SQLite & 34/162 & 106/162 \\
%       PostgreSQL & 0/162 & 11/162 \\
%       DuckDB & 0/162 & 13/162 \\
%       \bottomrule
%     \end{tabular}
%   \end{table}



% \para{Results}
% Table~\ref{tab:finetuned-rdbms-results} reports results\footnote{Tuning results for the 162-circuit dataset can be found here: \url{https://anonymous.4open.science/r/InferQ-ADEF/analysis/finetuned/rdbms162_all_engines_120s/rdbms162_all_engines_120s_results.csv}}
% for the tuned configurations on the 162-circuit validation set. Overall, Qiskit Aer
% remains the runtime winner on most circuits. The performance of PostgreSQL and DuckDB has
% improved substantially, but they do not achieve any runtime wins against the best
% successful Aer method on this subset. SQLite runtime does not improve significantly
% either. We note that for many of these circuits, Qiskit Aer completes each simulation in
% milliseconds, so further gains may require optimizing the end-to-end simulation workload rather than tuning each query separately.

% The primary benefit of tuning is improved memory efficiency: tuned SQLite wins peak memory
% on 106/162 circuits (up from 94 before tuning), and PostgreSQL and DuckDB improve from 0
% to 11/162 and 13/162 (both 0 before tuning) peak-memory wins, respectively.

% Over the large dataset of \num{7705} circuits, the results are similar. SQLite remains to be over fastest engine among RDBMSs. Our tuning did ot change much of its perfercentage of winning in runtime. But for memory, before tuning SQLite wins 3,196 circuits (41.6\%), after tuning 3,901. Umbra also wins 1 circuit. In totoal, RDBMS wins 3,902 50.6\%, around half of case. This makes RDBMS a promising simulation tool in terms of memory.

% \para{Take-away} Query optimization is foundamental but hard problem. I believe \sys shows quantum circuit simulation can be a new type of SQL query workload to optimze
% From these resutls together, we can see there is much things to explore, especially runtime. Combing the results from Appendix~\ref{app:qft-cap-sensitivity}, we think it is intersting to explore 1) optimize queries for single large circuit that requires performance degrade gracefully out of core; and 2) in quantum experiments they often run multiple simulation or same simulation repeatly. Then how to optimze a workload of multiple SQL queries representing such a simulation workload. \sys is a convient tool for generating SQL queries for such future research.
     


% \begin{table}[t]
%   \centering
%   \caption{SQLite static-contraction validation on the 7{,}705-circuit dataset. Winner counts compare tuned SQLite, untuned baseline SQLite, and the best successful Qiskit Aer method.}
%   \label{tab:sqlite-7705-static-results}
%   \footnotesize
%   \setlength{\tabcolsep}{3pt}
%   \begin{tabular}{lrrr}
%     \toprule
%     \textbf{Objective} & \textbf{Tuned SQLite} & \textbf{Baseline SQLite} & \textbf{Qiskit Aer} \\
%     \midrule
%     Runtime & 0 (0.0\%) & 1{,}067 (13.8\%) & 6{,}638 (86.2\%) \\
%     Peak memory & 2{,}085 (27.1\%) & 1{,}818 (23.6\%) & 3{,}802 (49.3\%) \\
%     \bottomrule
%   \end{tabular}
% \end{table}

% For the larger 7{,}705-circuit evaluation, Table~\ref{tab:sqlite-7705-static-results}
% shows the same pattern at scale.\footnote{\url{https://anonymous.4open.science/r/InferQ-ADEF/analysis/training_data/rdbms_training_data.parquet}}
% Qiskit Aer remains the runtime winner on 6{,}638/7{,}705 circuits (86.2\%). In contrast,
% SQLite is competitive on peak memory: tuned SQLite is the peak-memory winner on
% 2{,}085/7{,}705 circuits, and the baseline SQLite configuration wins on 1{,}818/7{,}705
% circuits. These results confirm that SQL-based simulation is not universally faster, but
% its sparse relational representation can be substantially more memory-efficient, which
% motivates learned routing and DBMS-level optimization for circuit-dependent performance.



\section{Indexes, partitioning, and CTE query structure}
\medskip
 
In \sys, each circuit is compiled into a \emph{single SQL statement} of the form
\texttt{WITH ... SELECT}, where each tensor contraction step is represented as a Common Table Expressions (CTE) and
implemented as a join followed by an aggregate, consistent with prior SQL-based
formulations \cite{hai2025quantum, einsteinSQL2023}. This CTE-chain representation is
useful because it makes the contraction order explicit and exposes the workload to the
DBMS optimizer as a sequence of standard relational operators (joins and group-bys).
However, it also limits classic physical design actions: intermediate CTE results are
 not persistent relations, so users cannot directly create secondary indexes, define
table partitioning, or maintain materialized views over those intermediates. Instead, the
engine may build transient internal structures (e.g., SQLite's automatic indexes and temporary B-tree structures), but these are optimizer-managed and not user-controlled. To enable explicit
physical design, we added a second query generator which lowers the same contraction pipeline
into a sequence of \texttt{CREATE TEMP TABLE AS} statements; each intermediate becomes a
first-class relation on which indexes and partitioning can be applied. 
Note that this is not part of \sys's standard pipeline; we added it to support the
experiments in this section. Systematically exploring index and partition selection for
these intermediate relations is a natural direction for future work.
%

 

\medskip
\para{Implications for physical design}
The tuned profiles in Table~\ref{tab:finetuned-rdbms-profiles} provide a reproducible
baseline that isolates DBMS-level execution effects (memory budgeting, temporary-workspace
management, parallel execution, and planner constraints) for this workload. Combined with
\texttt{split} execution, they also define a concrete path for future physical design
studies using standard database methodology. For example, one can treat intermediate
tables as design objects and evaluate: (i) index selection on join keys for contraction
steps (e.g., B-tree or hash indexes depending on the engine and operator), (ii)
partitioning and layout decisions to improve join locality and reduce spill (where the engine
supports partitioned tables and partitionwise planning), and (iii) selective
materialization and caching of frequently reused intermediates (e.g., repeated subcircuits)
as materialized relations. Because \sys exposes circuit features and query-shape metadata
(e.g., number of contraction steps (CTEs), runtime, and spill proxies), these choices can be
studied systematically as physical design problems across engines and workload regimes.

\medskip
\para{Preliminary Results}
We have conducted an experiment similar to Appendix~\ref{app:qft-cap-sensitivity} under a
4~GB memory limit. We choose a dense circuit from existing benchmarks, namely QFT, and vary
its qubit count from 3 to 26. Except for DuckDB at 26 qubits, all simulations completed.
We observe that the new query generator, which forces CTE materialization, reduces runtime
by roughly $2.1\times$ for DuckDB and $1.8\times$ for PostgreSQL; SQLite is essentially
unchanged. This shows that CTE materialization can be a potentially effective approach
for improving DBMS-backed simulation on dense workloads under tight memory limits.

% \medskip
% \para{Summary}
% Overall, this DBMS-specific track complements our untuned evaluation by (i) quantifying
% how sensitive SQL-based simulation is to engine configuration, and (ii) separating
% single-statement CTE execution from split execution, which is the appropriate regime for
% explicit indexes, partitioning, and materialized intermediates. The resulting tuned
% profiles and runners provide a reproducible baseline for both performance evaluation and
% future physical-design research on database-native circuit simulation.

% \para{Indexes, partition, and CTEs} In this work, for each circuit generated by \sys,  
% we compile it into a  \texttt{WITH ... SELECT} statement whose contraction
% steps are represented by a single CTE chain in the SQL query following existing work \cite{hai2025quantum, einsteinSQL2023}. 
% % Now we see that this query structure limits the opportunities for 
% % user-managed physical design: there are no durable intermediate relations on which one
% % can create explicit indexes or partitions. Instead, it relies on each
% % engine's internal choices, e.g., SQLite's automatic indexes and temporary B-tree
% % structures. In contrast, \texttt{split} execution lowers the same contraction plan into
% % a sequence of \texttt{CREATE TEMP TABLE AS} statements, making each intermediate an
% % explicit physical relation. This is the natural regime to study user-managed indexes,
% % materialized intermediates, and (where supported) partitioning planning.
% This matters for physical design. In this query
% shape there are no durable intermediate relations on which \sys can build
% user-managed indexes. Explicit indexes are therefore only applicable to split
% query mode, where the same CTE chain is lowered into per-step
% \texttt{CREATE TEMP TABLE} statements and each intermediate can be treated as a
% physical relation. In the monolithic profile we instead rely on each engine's
% internal planning choices, including automatic or ephemeral index structures
% where available, such as SQLite's automatic indexes. Partitioning has the same
% scope: PostgreSQL partitionwise planner switches can be enabled, but the
% monolithic CTEs are not declared partitioned tables. Partitioning is therefore a
% future split-mode or materialized-table design choice rather than a central knob
% in the two-profile monolithic experiment.

% The fine-tuned runner uses the raw monolithic IQS query shape: each circuit is
% compiled into a single \texttt{WITH ... SELECT} statement whose contraction
% steps are represented as CTEs. This matters for physical design. In this query
% shape there are no durable intermediate relations on which InferQ can build
% user-managed indexes. Explicit indexes are therefore only applicable to split
% query mode, where the same CTE chain is lowered into per-step
% \texttt{CREATE TEMP TABLE} statements and each intermediate can be treated as a
% physical relation. In the monolithic profile we instead rely on each engine's
% internal planning choices, including automatic or ephemeral index structures
% where available, such as SQLite's automatic indexes. Partitioning has the same
% scope: PostgreSQL partitionwise planner switches can be enabled, but the
% monolithic CTEs are not declared partitioned tables. Partitioning is therefore a
% future split-mode or materialized-table design choice rather than a central knob
% in the two-profile monolithic experiment.


 

\section{Additional Configuration Details on Experiments for Reproducibility}
\label{app:tuning-repro}

This section documents how we measure spilling performance and how we execute tuned runs.

\subsection{Circuit Construction Details}
\label{app:expander-details}
We design a circuit family with large qubit
counts but fixed sparse values.
Each circuit is parameterized by $(n,h)$, where $n$
is the number of qubits and $h<n$ is a small seed size. The construction applies Hadamards
only to the $h$ seed qubits and then applies only CNOT layers, which preserve support
size; therefore the final state has exactly $2^h$ non-zero amplitudes in the computational
basis. We
evaluate a sample of Expander circuits with $n \in [40,45]$, where a dense statevector
would require approximately 16--512~TB.  
An Expander circuit on $n$ qubits is parameterized by a seed size $h<n$. We apply
Hadamards to the first $h$ qubits to create a superposition over $2^h$ basis states and
initialize the remaining qubits to $\ket{0}$. We then apply only CNOT layers, which are
reversible linear transformations and preserve the number of basis states in support.
Consequently, the final state remains a uniform superposition over an affine subspace of
dimension $h$ and has exactly $2^h$ non-zero amplitudes in the computational basis (each
with magnitude $2^{-h/2}$). Figure~\ref{fig:expander_circuit} illustrates the circuit
structure.

\begin{figure}[t]
    \centering
    \includegraphics[width=1.0\linewidth]{images/expander_circuit_diagram.pdf}
    \caption{Expander circuit structure: Hadamards on $h$ seed qubits followed by CNOT-only
    layers that expand connectivity while preserving sparse support.}
    \label{fig:expander_circuit}
\end{figure}

Such sparse or highly structured states occur in practice in affine and stabilizer-style subroutines and
in workloads where computation is restricted to a low-dimensional subspace. Moreover,
even when the final state is dense, intermediate tensors and relations during contraction
can be much smaller than $2^n$. Separately, amplitude-amplification procedures (e.g.,
Grover-style iterations) tend to \emph{concentrate probability mass} onto a small subset
of basis states as they iterate, and sparse Hamiltonians and sparse linear systems (as in
eigensolver and HHL-style routines) often induce structured tensors. These regimes
motivate studying DBMS backends with sparse representations and out-of-core execution.

\subsection{Spill Configuration for Out-of-core experiments}
\label{app:spill-instrumentation}

We run SQL queries on PostgreSQL, DuckDB, and SQLite, sweeping Docker memory limit, e.g.,  16~GB. 
All containers use a one-CPU quota, swap is disabled by setting Docker
\texttt{--memory-swap} equal to the memory cap, and host page cache is not
explicitly dropped between runs. DuckDB uses one thread and sets
\texttt{memory\_limit} to the cap minus a 1024~MB headroom pad (3072, 7168,
and 15360~MB for the 4, 8, and 16~GB caps, respectively). SQLite uses a
64~MB cache budget for both main and temporary schemas. PostgreSQL uses the
tuned PostgreSQL~12.22 container, with the server container capped at the
experiment memory limit and \texttt{temp\_file\_limit} set to 64~GB. Spill is
measured as \texttt{postgres\_explain\_temp\_written} for PostgreSQL and as
\texttt{proc\_io\_write\_bytes} for DuckDB and SQLite.

We report spill volume using engine-specific sources when available, and otherwise use a
conservative operating system level proxy.

\noindent\textbf{PostgreSQL.} We compute exact temporary spill bytes from \texttt{EXPLAIN} (ANALYZE, BUFFERS, FORMAT JSON) by converting the reported temp written/read blocks to bytes using the 8192-byte block size; we treat this value as the per-query externalization metric.

\noindent\textbf{DuckDB.} We use \texttt{proc\_\allowbreak io\_\allowbreak write\_\allowbreak bytes} as the primary spill/workspace proxy, and when available we additionally parse DuckDB profiling JSON fields such as \texttt{temporary\_\allowbreak storage\_\allowbreak bytes} and \texttt{spilled\_\allowbreak bytes} for diagnostics only, since the profile schema can vary across versions.

\noindent\textbf{SQLite.} Similar to DuckDB, here we use \texttt{proc\_\allowbreak io\_\allowbreak write\_\allowbreak bytes} as the spill/workspace proxy. SQLite does not expose a stable query-level temp-byte counter through Python’s \texttt{sqlite3} interface, so we intentionally do not report a DB-native temp metric.

% \noindent\textbf{Qiskit Aer.} Aer does not spill, so we report no spill metric and instead record the simulator method, wall time, status, cgroup peak RSS, and whether execution fails due to memory.

\subsection{Configuration for DB tuning}
\label{app:mlos-runner}
We add a dedicated runner under \texttt{scripts/finetuned\_rdbms} that reuses the same
InferQ circuits and SQL generation path, but executes the emitted SQL under
engine-specific configurations rather than untuned defaults. The same directory contains
an automated tuning script (\texttt{mlos\_tune\_rdbms.py}) based on MLOS.

MLOS keeps the workload fixed (each circuit is compiled once into a SQL
query) and searches only over DBMS configuration parameters. For each engine, the tuner
evaluates seed profiles (e.g., \texttt{balanced}, \texttt{fast}, \texttt{spill\_safe})
and then uses a black-box optimizer backend (e.g., FLAML or SMAC) to propose additional
configurations. Its objective is the median runtime on a pilot set; timeouts and
execution failures receive a large penalty. All trials are recorded in
\texttt{mlos\_trials.csv}, and the best configuration per engine is stored in
\texttt{mlos\_best\_configs.json}. These winning configurations are then validated on the
full discovered circuit set via the fine-tuned runner using \texttt{--tuning-json}. Thus,
the reported improvements correspond to \emph{DBMS-level physical execution choices}
rather than changes to the circuit generator or SQL compilation.




Operationally, the runner generates SQL once per circuit, records query-shape metadata
(e.g., number of CTEs/contraction steps), and then executes the workload under the
selected engine profiles with warm-up and timed runs, configurable timeouts, and resume
support. To reduce total tuning cost, we optionally use SQLite as a lightweight gate:
if SQLite times out for a circuit under the tuned profile, we mark other engines as
\texttt{skipped\_sqlite\_timeout} for that circuit (i.e., we do not spend additional
compute on circuits already rejected by the gate). This produces a controlled comparison
of tuned physical execution while keeping the workload and SQL generation fixed.

% \para{Rationale for the two-profile settings.}
Table~\ref{tab:finetuned-rdbms-profiles} reports the best configurations selected and validated by our tuning workflow. For DuckDB, \texttt{small\_size} allocates a larger memory budget and modest parallelism for short query plans, while disabling insertion-order preservation to enable spill-capable operator choices; \texttt{large\_size} reduces memory and parallelism to limit concurrent memory pressure and temp-space exposure for deeper plans where planner/executor overhead can dominate. For SQLite, \texttt{small\_size} keeps the database and temporaries fully in-memory to avoid filesystem overhead, whereas \texttt{large\_size} uses a file-backed database with a modest mmap window to stabilize larger intermediate results while still keeping temporaries in memory. For PostgreSQL, \texttt{small\_size} increases per-node memory (\texttt{work\_mem}, \texttt{temp\_buffers}) and enables parallel workers to accelerate small contractions, while constraining join enumeration (\texttt{join\_collapse\_limit}/\texttt{from\_collapse\_limit}) to preserve the emitted join order and disabling JIT; \texttt{large\_size} caps per-node memory and disables parallel gather to reduce worker amplification and make externalization behavior more predictable. Finally, we include Umbra via its PostgreSQL-compatible frontend to broaden the engine comparison to a modern compiled analytical DBMS without changing the SQL workload.

\section{Generated Circuit Characterization}
\label{app:circuit-characterization}

\begin{figure*}[t]
    \centering
    \begin{subfigure}[t]{0.48\textwidth}
        \centering
        \includegraphics[width=\linewidth]{images/01_gate_mix.png}
        \caption{Top-20 gate-type distribution.}
        \label{fig:gate-mix}
    \end{subfigure}\hfill
    \begin{subfigure}[t]{0.48\textwidth}
        \centering
        \includegraphics[width=\linewidth]{images/03_depth_vs_width.png}
        \caption{Depth/width coverage.}
        \label{fig:depth-width}
    \end{subfigure}
    \caption{Generated-corpus gate mix and depth/width coverage. Color in \subref{fig:depth-width} indicates the number of circuits in each bin on a logarithmic scale.}
    \label{fig:generated-characterization-overview}
\end{figure*}

\begin{figure*}[t]
    \centering
    \includegraphics[width=\textwidth]{images/02_circuit_structure.png}
    \caption{Generated-corpus static circuit-feature distributions. Dashed lines mark medians.}
    \label{fig:circuit-structure}
\end{figure*}

\begin{figure*}[t]
    \centering
    \includegraphics[width=\textwidth]{images/04_graph_features.png}
    \caption{Generated-corpus interaction-graph feature distributions. Dashed lines mark medians.}
    \label{fig:interaction-graph}
\end{figure*}


To make the generated workload distribution transparent and reproducible, \sys records per-circuit metadata for each generated circuit and aggregates these fields over the full corpus. Figures~\ref{fig:generated-characterization-overview}, \ref{fig:circuit-structure}, and~\ref{fig:interaction-graph} summarize the resulting 202{,}975-circuit inventory from complementary perspectives: gate mix, static circuit structure, width/depth coverage, and interaction-graph structure. Together, these summaries show that the generator produces a broad workload rather than a collection of circuits that differ only in size.

\medskip
\para{Gate-mix distribution}
Figure~\ref{fig:gate-mix} reports the top-20 gate types by share of total gate count. The corpus has a substantial entangling component: \texttt{cx} alone accounts for 26.6\% of all gates, alongside common one-qubit rotations and Clifford-style operations such as \texttt{ry}, \texttt{x}, \texttt{h}, and \texttt{rz}. The remaining top gates include controlled-phase, controlled-Z, multi-controlled-X, swap, and higher-controlled variants, indicating that the generated workload includes both simple local transformations and larger multi-qubit control patterns.


\medskip
\para{Static circuit structure}
Figure~\ref{fig:circuit-structure} expands the view from gate identities to circuit-level structure. The generated corpus spans small and moderate-width circuits, with median values of 14 qubits, width 16, depth 59, and circuit size 190. It also varies substantially in two-qubit-gate count and two-qubit-gate percentage, which are important because they influence both tensor contraction structure and the amount of intermediate state growth during simulation. Additional static descriptors such as Pauli-gate count, locality ratio, idling score, and density score capture whether a circuit is mostly local or contains broader qubit interactions and idle regions. These distributions make the benchmark more informative than a width/depth grid alone.

\medskip
\para{Depth/width coverage}
Figure~\ref{fig:depth-width} shows the joint coverage over number of qubits and circuit depth. The heat map reports the number of circuits in each \emph{(number of qubits, depth)} bin using a logarithmic color scale, exposing both dense regions of the generated space and sparsely populated edge cases. The workload includes shallow circuits, deep circuits with depths above $10^3$, and circuit sizes extending to roughly 30 qubits in this corpus. This joint view is useful for interpreting performance results because simulator and DBMS behavior can change sharply when both width and depth increase, even when either feature alone appears moderate.

\medskip
\para{Interaction-graph structure}
Figure~\ref{fig:interaction-graph} characterizes the qubit interaction graphs induced by the generated circuits. We construct an interaction graph by treating qubits as vertices and adding an edge when two qubits participate in a multi-qubit operation. The resulting features include edge count, maximum and average degree, clustering coefficient, shortest-path statistics, diameter, radius, cut-related measures, central-point dominance, and adjacency-matrix variation. These graph features capture structure that is not visible from gate count, width, or depth alone: two circuits with the same number of qubits and a similar depth can still differ substantially in interaction density, path length, graph diameter, or degree distribution. Such differences affect SQL query shape, contraction opportunities, intermediate-result sizes, and simulator behavior.

Overall, these summaries make the generated workload transparent, reproducible, and analytically useful. With \sys, users can inspect the concrete properties that shape simulation performance, including gate composition, static circuit size, depth/width coverage, and qubit-interaction topology. This allows engine comparisons to be interpreted in terms of the circuits being simulated and helps identify which circuit families require particular simulation strategies.